\documentclass[11pt]{article}

\usepackage[margin=1in]{geometry}
\usepackage{graphicx}
\usepackage{array}
\usepackage{hyperref}
\usepackage{longtable}
\usepackage{enumitem}
\usepackage{booktabs}
\usepackage{caption}
\usepackage{amsmath}
\usepackage{titlesec}
\usepackage{float}

\titleformat{\section}{\large\bfseries}{\thesection.}{1em}{}

\begin{document}

\begin{center}
    \Large \textbf{Sri Sivasubramaniya Nadar College of Engineering, Chennai} \\
    \large (An Autonomous Institution affiliated to Anna University) \\
    \vspace{0.3cm}
\end{center}

\begin{table}[h!]
\renewcommand{\arraystretch}{1.5}
\centering
\resizebox{\textwidth}{!}{%
\begin{tabular}{|l|cll|}
\hline
\textbf{Degree \& Branch} & \multicolumn{1}{c|}{B.E Computer Science \& Engineering} & \textbf{Semester} & VI \\ \hline
\textbf{Subject Code \& Name} & \multicolumn{3}{c|}{UCS2612 -- Machine Learning Laboratory} \\ \hline
\textbf{Academic Year} & \multicolumn{1}{c|}{2025--2026 (Even)} & \textbf{Batch} & 2023--2027 \\ \hline
\textbf{Name} & \multicolumn{1}{c|}{Mehanth T} & \textbf{Register No.} & 3122235001080 \\ \hline
\end{tabular}
}
\end{table}

\begin{center}
    \textbf{\Large Experiment 5}
\end{center}

\begin{center}
    \textbf{Perceptron vs Multilayer Perceptron with Hyperparameter Tuning}
\end{center}

\section*{Objective}
To implement and compare the performance of:
\begin{itemize}
    \item \textbf{Model A:} Single-Layer Perceptron Learning Algorithm (PLA)
    \item \textbf{Model B:} Multilayer Perceptron (MLP) with hidden layers and nonlinear activation functions
\end{itemize}

Students must select, tune, and justify hyperparameters such as learning rate, activation function, optimizer, batch size, and number of hidden layers.

\section*{Dataset}
\textbf{English Handwritten Characters Dataset}

\textbf{Download Link:}  
\href{https://www.kaggle.com/datasets/dhruvildave/english-handwritten-characters-dataset}{https://www.kaggle.com/datasets/dhruvildave/english-handwritten-characters-dataset}

\begin{itemize}
    \item Total samples: 3,410
    \item Number of classes: 62 (0--9, A--Z, a--z)
    \item Image type: Grayscale
\end{itemize}

\section*{Theory}

\subsection*{Perceptron Learning Algorithm (PLA)}
The Perceptron is a linear classifier that uses a step activation function to separate data into two classes.

\[
w_{t+1} = w_t + \eta (y - \hat{y})x
\]

\begin{itemize}
    \item Uses a step activation function
    \item Converges only for linearly separable data
    \item Extended to multi-class problems using One-vs-Rest strategy
\end{itemize}

\textbf{Limitation:} PLA cannot learn nonlinear decision boundaries required for handwritten character recognition.

\subsection*{Multilayer Perceptron (MLP)}
An MLP consists of one or more hidden layers with nonlinear activation functions.

\begin{itemize}
    \item Activation functions: ReLU, Sigmoid, Tanh
    \item Loss function: Categorical Cross-Entropy
    \item Optimizers: SGD, Adam
    \item Learns nonlinear decision boundaries using backpropagation
\end{itemize}

\section*{Steps for Implementation}

\begin{enumerate}[label=\arabic*.]
    \item Preprocess the dataset (resize, flatten, normalize, \textbf{StandardScaler}).
    \item Implement \textbf{PLA} from scratch:
    \begin{itemize}
        \item Use step activation function.
        \item Apply perceptron weight update rule.
        \item Extend PLA to multi-class classification using One-vs-Rest.
    \end{itemize}
    \item Implement \textbf{MLP}:
    \begin{itemize}
        \item Select number of hidden layers and neurons.
        \item Choose activation functions.
        \item Train using backpropagation.
    \end{itemize}
    \item Perform \textbf{hyperparameter tuning}:
    \begin{itemize}
        \item Learning rate
        \item Batch size
        \item Optimizer
        \item Activation function
    \end{itemize}
    \item Compare PLA and MLP using evaluation metrics.
\end{enumerate}

\section*{Performance Evaluation}

\begin{table}[h!]
\centering
\caption{Overall Performance Comparison between PLA and MLP}
\renewcommand{\arraystretch}{1.3}
\begin{tabular}{lcccc}
\toprule
\textbf{Model} & \textbf{Accuracy (\%)} & \textbf{Precision} & \textbf{Recall} & \textbf{F1-score} \\
\midrule
PLA (OvR) & 9.24\% & 0.09 & 0.09 & 0.09 \\
MLP (Tuned) & 41.50\% & 0.43 & 0.41 & 0.41 \\
\bottomrule
\end{tabular}
\end{table}

\begin{table}[h!]
\centering
\caption{Hyperparameter Tuning Results for MLP}
\renewcommand{\arraystretch}{1.3}
\begin{tabular}{cccccc}
\toprule
\textbf{Hidden Layers} & \textbf{Activation} & \textbf{Optimizer} & \textbf{Learning Rate} & \textbf{Batch Size} & \textbf{Accuracy (\%)} \\
\midrule
1 (128) & ReLU & SGD & 0.01 & 32 & < 30\% \\
2 (256--128) & ReLU & Adam & 0.001 & 64 & \textbf{41.50\%} \\
3 (512--256--128) & Tanh & Adam & 0.0005 & 64 & < 30\% \\
\bottomrule
\end{tabular}
\end{table}

\begin{table}[h!]
\centering
\caption{Training Convergence Comparison}
\renewcommand{\arraystretch}{1.3}
\begin{tabular}{lccc}
\toprule
\textbf{Model} & \textbf{Epochs} & \textbf{Final Training Loss} & \textbf{Convergence Behavior} \\
\midrule
PLA & 10 & N/A & Poor Convergence \\
MLP (Tuned) & 300 & Low & Stable convergence (Early Stopping) \\
\bottomrule
\end{tabular}
\end{table}

\section*{Evaluation Metrics}
The MLP model significantly outperformed the PLA model. The confusion matrix and loss curves (generated in plots/) demonstrate that the MLP effectively learned features, whereas the PLA struggled with the non-linear complexity of the image data. The critical factor for MLP success was the application of \textbf{StandardScaler} to normalize input features to zero mean and unit variance.

\section*{A/B Experiment Comparison}
\begin{itemize}
    \item \textbf{Final Tuned Hyperparameters:} Hidden Layers: (256, 128), Activation: ReLU, Optimizer: Adam, Learning Rate: 0.001, Batch Size: 64.
    \item \textbf{Strengths/Weaknesses:} PLA is simple but fails on complex image data (9\% accuracy). MLP is computationally intensive but achieves much higher accuracy (41.5\%) by capturing non-linear relationships.
    \item \textbf{Tuning Effect:} The 'Adam' optimizer and 2-layer architecture provided the best stability and performance compared to SGD and single-layer configurations.
\end{itemize}

\section*{Observation Questions}
\begin{itemize}
    \item \textbf{Why does PLA underperform compared to MLP?} PLA is a linear classifier and cannot capture the complex, non-linear pixel relationships inherent in handwritten characters. MLP uses non-linear activation functions (ReLU) to model these boundaries.
    \item \textbf{Which hyperparameter had the most impact on MLP performance?} Feature Scaling (StandardScaler) was the most critical preprocessing step. Architecturally, the number of hidden layers and the choice of 'Adam' optimizer were significant.
    \item \textbf{How does optimizer choice affect convergence?} Adam (Adaptive Moment Estimation) generally converges faster and more reliably on sparse/complex data like images compared to standard SGD, which can get stuck in local minima or oscillate.
    \item \textbf{Does increasing hidden layers always improve performance?} No. While more layers add capacity, they also increase the risk of overfitting and training difficulty (vanishing/exploding gradients). In this experiment, 2 layers worked better than 3 or 1.
    \item \textbf{How can overfitting be detected and mitigated?} Overfitting is detected when training accuracy is high but validation/test accuracy is low. It can be mitigated using regularization (L2), Dropout, or Early Stopping (used in this experiment).
\end{itemize}

\end{document}
